% LTeX: enabled=false

% Here you can see an example of how to create text conditioned by the language
% variable. The \iftoggle command:
%
%   \iftoggle{ThesisInEnglish}{%
%   <your-text-in-english>
%   }{%
%   <your-text-in-french>
%   }
%
% will compile only one of the two blocks, depending on the variable you set at
% the beginning of this document. Language selection is managed this way in the
% formatAndDefs.tex file. You too can create sections of your thesis that is
% language dependend this way, although you probably won't need it. Another use
% of \iftoggle can be found at the end of this file.
\iftoggle{ThesisInEnglish}{%
\section*{Acknowledgments}
}{%
\section*{Remerciements}
}

% Un tel travail aurait été impossible à réaliser sans la gentillesse, l'aide et le soutien de celles et ceux qui m'ont aidé.\\

% Mes remerciements vont d'abord à Paolo ROBUFFO GIORDANO et Pascal MORIN, qui ont accepté de relire cette thèse et d'en être rapporteur. La version définitive de ce mémoire a bénéficié de leur lecture très attentive et de leurs remarques précieuses. Je tiens à remercier Sophie TARBOURIECH d’avoir accepté d’être présidente du jury. Je remercie également Philippe CHEVREL d’avoir accepté d’assister à la présentation de ce travail.\\

% Je tiens à remercier l'ensemble de mes encadrants, Fabrice DEMOURANT, Gautier HATTENBERGER, Thomas LOQUEN et Luca ZACCARIAN pour leurs bienveillances, leurs conseils avisés et les nombreuses heures qu'ils m'ont données. \\

% J'ai une attention particulière pour Luca, tu m'as accompagné depuis le début de mon stage de master jusqu'à la fin de mon doctorat. J'ai énormément appris à tes côtés et je te remercie pour toute l'énergie que tu as mise dans ma réussite.\\

% Gautier, j'ai pu appréhender un système de gestion de drone complexe grâce à toi. J'observe aujourd'hui le résultat de l'ensemble des efforts que j'ai mis pour devenir autonome. Je ne peux que te remercier pour ces nombreuses explications, pour le temps que tu as passé à m'expliquer et pour les nombreux conseils.\\

% Merci à Fabrice pour ces nombreux échanges autour de la commande robuste. Je suis fier du travail que nous avons mené ensemble.\\

% Thomas, je te remercie d'avoir été à mes côtés à chaque instant, prévenant et attentionné, tu as toujours fait en sorte que mon doctorat se déroule dans les meilleures conditions.\\

% Je n'aurais jamais obtenu de tels résultats sans une équipe aussi diverse et intéressante dans laquelle vous m'avez intégré. Merci à Murat BRONZ pour nos échanges autour des drones convertibles. Merci à Alexandre BUSTICO pour m'avoir octroyé du temps pour résoudre mes problèmes d'implémentation. Merci à Michel GORRAZ pour les nombreuses discussions autour de l'électronique, mais pas seulement. Enfin, merci à Fabien BONNEVAL pour les nombreuses heures que nous avons partagé à coder ou à discuter des drones.\\

% Merci à Jim pour les nombreux échanges.
% Merci à Thierry pour m'avoir permis d'enseigner aux seins de l'ENAC et pour nos nombreuses sorties course à pied.\\

% Enfin, je tiens à remercier tout le laboratoire drone pour les nombreux moments de partage, les repas, les pauses-café, les conférences imav et tous les moments qui ont été pour moi des motivations à finir mon travail.\\

% Je tiens à remercier ma famille d'avoir participé à ma soutenance. Ce moment fut l'occasion pour moi de vous présenter une partie de mon travail.\\


% Ma dernière pensée va à ma femme. Cette année 2024 fut une grande année, car nous nous sommes mariés et j'ai fini mon doctorat. Ce travail t'appartient autant qu'à moi, puisque c'est grâce à toi que je l'ai mené à son terme. Tu as toujours été là pour me soutenir dans les nombreux moments de doute, pour te réjouir de mes réussites et pour me pousser quand il était nécessaire. 


Un tel travail n’aurait jamais pu voir le jour sans la gentillesse, le soutien et l’aide précieuse de toutes celles et ceux qui m’ont accompagné.

Mes remerciements vont tout d’abord à Paolo Robuffo Giordano et Pascal Morin, qui ont accepté de relire cette thèse et d’en être rapporteurs. La version finale de ce mémoire a largement bénéficié de leur lecture attentive et de leurs remarques pertinentes. Je tiens également à remercier Sophie Tarbouriech pour avoir accepté la présidence du jury, ainsi que Philippe Chevrel pour sa présence et son intérêt lors de la présentation de ce travail.

Je suis profondément reconnaissant envers l’ensemble de mes encadrants, Fabrice Demourant, Gautier Hattenberger, Thomas Loquen et Luca Zaccarian, pour leur bienveillance, leurs conseils avisés et le temps précieux qu’ils m’ont accordé.

Une attention particulière va à Luca, qui m’a accompagné depuis le début de mon stage de master jusqu’à la fin de ce doctorat. J’ai énormément appris à tes côtés, et je te suis infiniment reconnaissant pour toute l’énergie que tu as consacrée à ma réussite.

Gautier, grâce à toi, j’ai pu appréhender un système complexe de gestion de drones. Aujourd’hui, je mesure pleinement les fruits des efforts investis pour devenir autonome. Je te remercie sincèrement pour ton temps, tes nombreuses explications et tes précieux conseils.

Merci à Fabrice pour nos nombreux échanges passionnants autour de la commande robuste. Je suis fier du travail accompli ensemble.

Thomas, tu as été un soutien constant et attentionné, veillant à ce que mon doctorat se déroule dans les meilleures conditions possibles. Je t’en suis profondément reconnaissant.

Je tiens également à remercier les membres de l’équipe qui m’ont entouré et enrichi :
\begin{itemize}
    \item Murat Bronz, pour nos échanges constructifs sur les drones convertibles,
    \item Alexandre Bustico, pour le temps que tu m’as accordé pour résoudre mes problèmes d’implémentation,
    \item Michel Gorraz, pour nos discussions enrichissantes autour de l’électronique et bien d’autres sujets encore,
    \item Fabien Bonneval, pour ces nombreuses heures passées ensemble à coder et à partager nos réflexions sur les drones.
\end{itemize}


Merci également à Jim pour nos nombreux échanges, et à Thierry pour m’avoir offert l’opportunité d’enseigner au sein de l’ENAC, ainsi que pour nos mémorables sorties de course à pied.

Enfin, un grand merci à tout le laboratoire drone pour les moments de partage, les repas, les pauses-café, les conférences IMAV et tous ces instants qui ont été pour moi de véritables sources de motivation pour mener ce travail à son terme.

Je remercie chaleureusement ma famille pour sa présence lors de ma soutenance, un moment unique où j’ai eu la chance de vous présenter une partie de mon travail.

Ma dernière pensée va à ma femme. Cette année 2024 restera inoubliable : celle de notre mariage et de l’aboutissement de mon doctorat. Ce travail t’appartient autant qu’à moi, car c’est grâce à toi que j’ai pu le mener jusqu’au bout. Tu as toujours été présente dans les moments de doute, tu t’es réjouie de mes réussites et tu m’as poussé lorsque cela était nécessaire. Pour tout cela, je te suis éternellement reconnaissant.